\documentclass[11pt]{article}
\usepackage{fontspec}
\setmainfont{DejaVu Serif}
\usepackage{amsmath,amssymb}
\usepackage{geometry}
\geometry{margin=2.4cm}
\usepackage{xcolor}
\usepackage{graphicx}
\usepackage{enumitem}
\usepackage[hidelinks]{hyperref}

\definecolor{ink}{HTML}{1a2a3a}
\definecolor{accent}{HTML}{1f5fa8}
\definecolor{soft}{HTML}{eef3fb}
\color{ink}

\newcommand{\key}[1]{\par\medskip\noindent\colorbox{soft}{\parbox{\dimexpr\linewidth-2\fboxsep}{\textcolor{accent}{\textbf{Κλειδί:}}\ #1}}\par\medskip}
\newcommand{\ex}[1]{\par\smallskip\noindent\textbf{\textcolor{accent}{Παράδειγμα.}}\ \textit{#1}\par\smallskip}

\setlength{\parskip}{0.5em}
\setlength{\parindent}{0pt}

\title{\textbf{Φροντιστήριο για αρχάριους}\\[2pt]
\large Τα μαθηματικά πίσω από το PEAD event-study\\[-2pt]
\normalsize\textcolor{accent}{Πώς μετράμε αν τα κέρδη προβλέπουν την τιμή — και για πόσο}}
\author{QuantIQ / AGEL AI}
\date{}

\begin{document}
\maketitle
\thispagestyle{empty}

\noindent Αυτό το φυλλάδιο εξηγεί, \emph{από το μηδέν}, τα μαθηματικά που χρησιμοποιεί το module
\texttt{pead\_event\_study.py}. Δεν χρειάζεσαι προηγούμενες γνώσεις πέρα από βασική άλγεβρα.
Κάθε έννοια έρχεται με ένα μικρό αριθμητικό παράδειγμα. Στο τέλος θα καταλαβαίνεις
\emph{τι} μετράμε, \emph{γιατί}, και \emph{τι σημαίνει} κάθε αριθμός που βγάζει ο κώδικας.

\tableofcontents
\bigskip

%==================================================================
\section{Τι προσπαθούμε να μετρήσουμε (με λόγια)}

Όταν μια εταιρεία ανακοινώνει κέρδη \emph{καλύτερα από το αναμενόμενο}, η μετοχή της δεν
ανεβαίνει αμέσως όσο «πρέπει». Συνεχίζει να ανεβαίνει σιγά-σιγά για εβδομάδες. Αυτό λέγεται
\textbf{PEAD} (Post-Earnings-Announcement Drift — μετα-ανακοινωσιακή ολίσθηση). Οι επενδυτές
\emph{υπο-αντιδρούν} και «χωνεύουν» την είδηση σταδιακά.

Θέλουμε να απαντήσουμε τρία ερωτήματα με αριθμούς:
\begin{enumerate}[leftmargin=2em]
  \item Πόσο μεγάλη ήταν η «έκπληξη» στα κέρδη; \quad(\textbf{SUE})
  \item Πόσες μέρες κρατάει η ολίσθηση και πόσο γρήγορα σβήνει; \quad(\textbf{drift \& half-life})
  \item Είναι το φαινόμενο αληθινό, και αντέχει στον χρόνο; \quad(\textbf{στατιστική \& durability})
\end{enumerate}

%==================================================================
\section{Μέρος Α — Τα εργαλεία (γρήγορη επανάληψη)}

\subsection{Μέση τιμή}
Ο μέσος όρος $N$ αριθμών:
\[ \bar{x}=\frac{1}{N}\sum_{i=1}^{N}x_i . \]
\ex{Για $2,\ 4,\ 9$: $\ \bar{x}=(2+4+9)/3=5$.}

\subsection{Διασπορά και τυπική απόκλιση — «πόσο κουνιέται»}
Η \textbf{τυπική απόκλιση} $\sigma$ μετράει πόσο απλώνονται οι τιμές γύρω από τον μέσο:
\[ \sigma=\sqrt{\frac{1}{N}\sum_{i=1}^{N}(x_i-\bar{x})^2}. \]
Μικρό $\sigma$ = σταθερό· μεγάλο $\sigma$ = θορυβώδες.
\ex{Για $2,4,9$ με $\bar{x}=5$: αποκλίσεις $-3,-1,4$, τετράγωνα $9,1,16$, μέσο $26/3\approx8{,}67$,
άρα $\sigma=\sqrt{8{,}67}\approx2{,}9$.}

\subsection{Τυποποίηση (z-score) — κοινή «γλώσσα»}
Για να συγκρίνουμε μεγέθη με διαφορετικές κλίμακες, τα φέρνουμε σε κοινή μονάδα:
\[ z=\frac{x-\bar{x}}{\sigma}. \]
Ένα $z=+2$ σημαίνει «2 τυπικές αποκλίσεις πάνω από τον μέσο» — \emph{ασυνήθιστα} μεγάλο.

\subsection{Συσχέτιση — πάνε μαζί;}
Η συσχέτιση $\rho$ δύο μεγεθών $x,y$ παίρνει τιμές από $-1$ έως $+1$:
\[ \rho=\frac{\sum (x_i-\bar{x})(y_i-\bar{y})}{\sqrt{\sum(x_i-\bar{x})^2}\,\sqrt{\sum(y_i-\bar{y})^2}}. \]
$\rho=+1$: κινούνται απόλυτα μαζί· $\rho=0$: καμία σχέση· $\rho=-1$: αντίθετα.
Στις αγορές, ένα σήμα με $\rho\approx0{,}03$ με τις μελλοντικές αποδόσεις είναι \emph{αδύναμο αλλά
ίσως χρήσιμο}. Αυτή τη συσχέτιση τη λέμε \textbf{IC} (Information Coefficient).

\subsection{Κλίση παλινδρόμησης — πόσο αλλάζει το ένα ανά μονάδα του άλλου}
Αν θέλουμε «πόσο ανεβαίνει το $y$ όταν το $x$ ανεβαίνει κατά 1», βρίσκουμε την \textbf{κλίση}:
\[ b=\frac{\sum x_i\,y_i}{\sum x_i^{2}}\quad(\text{όταν τα }x\text{ έχουν μέσο 0}). \]
Είναι η ίδια ιδέα με μια ευθεία $y\approx b\,x$ που περνά «όσο καλύτερα γίνεται» από τα σημεία.

\subsection{Εκθετική φθορά και λογάριθμος — η «μισή ζωή»}
Πολλά πράγματα σβήνουν εκθετικά: ξεκινούν δυνατά και μειώνονται κατά σταθερό \emph{ποσοστό}:
\[ y(k)=b_0\,e^{-k/\tau}. \]
Το $\tau$ (ταυ) είναι ο «χρόνος φθοράς». Η \textbf{μισή ζωή} (half-life) είναι ο χρόνος για να
πέσει στο μισό:
\[ t_{1/2}=\tau\ln 2\approx 0{,}693\,\tau. \]
Με λογάριθμο, η εκθετική γίνεται ευθεία — γι' αυτό τη «ισιώνουμε» για να βρούμε το $\tau$:
\[ \ln y(k)=\ln b_0-\frac{1}{\tau}\,k. \]
\ex{Αν $\tau=15$ μέρες, τότε $t_{1/2}=15\times0{,}693\approx10{,}4$ μέρες: κάθε $\sim$10 μέρες
το σήμα μισαίνει.}

%==================================================================
\section{Μέρος Β — Από τα κέρδη στο σήμα}

\subsection{SUE: πόσο μεγάλη ήταν η έκπληξη}
Δεν χρειαζόμαστε προβλέψεις αναλυτών. Η κλασική μέθοδος (Foster 1977· Bernard \& Thomas 1989)
συγκρίνει τα κέρδη με το \emph{ίδιο τρίμηνο πέρυσι} (seasonal random walk):
\[ \boxed{\ \mathrm{SUE}_{i,t}=\frac{E_{i,t}-E_{i,t-4}}{\sigma_{i}}\ } \]
όπου $E_{i,t}$ τα κέρδη ανά μετοχή του τριμήνου, $E_{i,t-4}$ το ίδιο τρίμηνο πριν έναν χρόνο,
και $\sigma_i$ η τυπική απόκλιση αυτών των ετήσιων μεταβολών (τα τελευταία $\sim$8 τρίμηνα).

\ex{Φέτος $E_{t}=1{,}30$, πέρυσι $E_{t-4}=1{,}05$, άρα έκπληξη $=+0{,}25$.
Αν οι παλιές μεταβολές είχαν $\sigma=0{,}055$, τότε
$\mathrm{SUE}=0{,}25/0{,}055\approx +4{,}5$ — \emph{μεγάλη} θετική έκπληξη.}

\key{Το SUE δεν το κρίνουμε απόλυτα, αλλά \emph{συγκριτικά} με τις άλλες μετοχές την ίδια μέρα
(το τυποποιούμε με z-score). Long τα υψηλά SUE, short τα χαμηλά.}

\subsection{«Abnormal» αποδόσεις: αφαιρούμε την αγορά}
Αν όλη η αγορά ανέβηκε 1\% μια μέρα, δεν είναι «επίτευγμα» της μετοχής. Κρατάμε μόνο το
\emph{σχετικό} κομμάτι, αφαιρώντας τον μέσο όρο όλων των μετοχών εκείνης της μέρας:
\[ \mathrm{AR}_{i,t}=r_{i,t}-\bar{r}_{t},\qquad \bar{r}_t=\frac{1}{N}\sum_{j}r_{j,t}. \]
Έτσι το χαρτοφυλάκιο γίνεται \textbf{market-neutral} (ουδέτερο ως προς την αγορά): κερδίζει
μόνο αν διαλέγει \emph{σωστά μέσα} στην αγορά, όχι επειδή «ανέβηκαν όλα».

\subsection{Χρόνος γεγονότος (event time) και ο κανόνας T+1}
Ευθυγραμμίζουμε όλες τις ανακοινώσεις σε μια κοινή «μέρα 0». Δεν αγοράζουμε την ίδια μέρα της
ανακοίνωσης (δεν θα προλαβαίναμε ρεαλιστικά) — μπαίνουμε την \textbf{επόμενη} εργάσιμη
($T{+}1$). Μετά παρακολουθούμε τις μέρες $k=1,2,3,\dots$ μετά την είσοδο.

%==================================================================
\section{Μέρος Γ — Η καμπύλη ολίσθησης και φθοράς}

\subsection{Το «προφίλ ολίσθησης» $b(k)$}
Ρωτάμε: \emph{την ημέρα $k$ μετά την είσοδο, πόση απόδοση παίρνουμε ανά μονάδα SUE;}
Παίρνουμε όλα τα γεγονότα, τυποποιούμε το SUE ($z_e$), και υπολογίζουμε την κλίση:
\[ \boxed{\ b(k)=\frac{\sum_{e} z_e\,\mathrm{AR}_{e,k}}{\sum_{e} z_e^{2}}\ } \]
Είναι ακριβώς η «κλίση παλινδρόμησης» του Μέρους Α: πόση abnormal απόδοση φέρνει κάθε μονάδα
έκπληξης, την κάθε μέρα.

\ex{Τρία γεγονότα με $z=(+2,\,0,\,-1)$ και απόδοση τη μέρα $k$ ίση με
$\mathrm{AR}=(+0{,}008,\,+0{,}001,\,-0{,}003)$.
Αριθμητής $=2(0{,}008)+0(0{,}001)+(-1)(-0{,}003)=0{,}019$.
Παρονομαστής $=2^2+0+(-1)^2=5$.
Άρα $b(k)=0{,}019/5=0{,}0038=38$ μονάδες βάσης ανά μονάδα SUE.}

\subsection{Σωρευτική ολίσθηση $B(h)$ και βέλτιστος ορίζοντας}
Αν κρατήσουμε $h$ μέρες, η συνολική ολίσθηση ανά μονάδα SUE είναι το άθροισμα:
\[ B(h)=\sum_{k=1}^{h} b(k). \]
Η $B(h)$ ανεβαίνει όσο υπάρχει ολίσθηση και μετά \emph{επιπεδώνει} — εκεί που επιπεδώνει,
το drift εξαντλήθηκε, και αυτός είναι ο λογικός ορίζοντας κρατήματος.

\subsection{Φθορά και μισή ζωή}
Το $b(k)$ τυπικά φθίνει εκθετικά: $b(k)\approx b_0 e^{-k/\tau}$. Ισιώνοντάς το με λογάριθμο
(Μέρος Α) βρίσκουμε το $\tau$, και άρα τη μισή ζωή $t_{1/2}=\tau\ln 2$:
πόσες μέρες χρειάζονται για να πραγματοποιηθεί το μισό του drift.

\key{Στη δοκιμή μας (με \emph{γνωστή} αλήθεια), ο κώδικας ανέκτησε μισή ζωή
$\approx 8{,}6$ μέρες ενώ η πραγματική ήταν $10{,}4$ — πολύ κοντά. Δηλαδή η μέθοδος «βλέπει»
σωστά πόσο γρήγορα σβήνει το σήμα.}

\begin{center}
\includegraphics[width=\linewidth]{pead_results.png}\\[2pt]
{\small Αριστερά: το μετρημένο $b(k)$ (μπάρες) και η εκθετική προσαρμογή (κόκκινη, $t_{1/2}\approx9$ μέρες).
Μέση: η σωρευτική $B(h)$. Δεξιά: η απόδοση ανά έτος — η «αντοχή» (durability).}
\end{center}

%==================================================================
\section{Μέρος Δ — Είναι αληθινό; (στατιστική)}

\subsection{Γιατί φτιάχνουμε «ημερολογιακό» χαρτοφυλάκιο}
Δύο προβλήματα κάνουν τα απλά τεστ να \emph{ψεύδονται}:
\begin{itemize}[leftmargin=1.6em]
  \item \textbf{Επικαλυπτόμενα παράθυρα:} κρατάμε κάθε στοίχημα 60 μέρες, οπότε οι αποδόσεις
        διαδοχικών ημερών «μοιράζονται» πληροφορία (αυτοσυσχετίζονται).
  \item \textbf{Συγκέντρωση στο earnings season:} πολλές ανακοινώσεις μαζεύονται τις ίδιες
        εβδομάδες, άρα τα στοιχήματα δεν είναι ανεξάρτητα.
\end{itemize}
Λύση: κάθε μέρα φτιάχνουμε ένα χαρτοφυλάκιο από \emph{όλα} τα ενεργά γεγονότα (long/short με
βάρος το $z_e$), και κρατάμε μία ημερήσια σειρά αποδόσεων. Η ανάλυση σε «ημερολογιακό χρόνο»
αντιμετωπίζει σωστά την επικάλυψη και τη συγκέντρωση.

\subsection{Newey--West: τίμια $t$-στατιστική}
Η συνηθισμένη $t$-στατιστική υποθέτει \emph{ανεξάρτητες} μέρες. Όταν οι μέρες
αυτοσυσχετίζονται, η κανονική $t$ είναι \emph{υπερβολικά αισιόδοξη} (δείχνει «σημαντικό» ενώ
δεν είναι). Η διόρθωση \textbf{Newey--West} μεγαλώνει σωστά το «σφάλμα» λαμβάνοντας υπόψη τη
συσχέτιση μέχρι $L$ μέρες πίσω:
\[ t=\frac{\bar{r}}{\mathrm{SE}_{\text{NW}}},\qquad
   \mathrm{SE}_{\text{NW}}\ \text{μεγαλύτερο όταν οι μέρες συσχετίζονται.} \]
Πρακτικός κανόνας: όσο μεγαλύτερο το $|t|$, τόσο πιο απίθανο να είναι τύχη. Στις αγορές, λόγω
του ότι δοκιμάζουμε πολλά, θέλουμε αυστηρό κατώφλι (συχνά $|t|>3$).

\subsection{Information Ratio (IR)}
Η απόδοση ανά μονάδα ρίσκου (σε ετήσια βάση):
\[ \mathrm{IR}=\sqrt{252}\;\frac{\overline{r}}{\sigma_r}, \]
όπου $252$ οι εργάσιμες μέρες του έτους. Ρεαλιστικά «καλό» είναι $\mathrm{IR}\approx0{,}5$--$1{,}0$·
πολύ μεγαλύτερο, υποψιάσου αυταπάτη.

%==================================================================
\section{Μέρος Ε — Πόσο κρατάει το σήμα;}

\subsection{Αντοχή στον χρόνο (durability)}
Υπολογίζουμε το IR \emph{ανά έτος} και κοιτάμε αν φθίνει. Προσαρμόζουμε μια ευθεία
$\mathrm{IR}(\text{έτος})\approx a+(\text{κλίση})\cdot\text{έτος}$: αρνητική κλίση σημαίνει ότι
το φαινόμενο \emph{σβήνει} με τα χρόνια (έχει «αρμπιτραριστεί»). Αυτό είναι το πιο σημαντικό
μέγεθος για προϊόν: ένα σήμα που φθίνει γρήγορα σε αναγκάζει σε ατελείωτο κυνήγι.

\key{Στη δοκιμή, ο κώδικας έπιασε καθαρά τη φθορά που είχαμε «κρύψει»: από IR $\sim$5 στα πρώτα
χρόνια σε $\sim$0 στα τελευταία (κλίση $\approx -0{,}30$ ανά έτος).}

\subsection{Επιμονή του SUE (το νόμιμο «forward» κομμάτι)}
Δεν μπορούμε να προβλέψουμε την \emph{επόμενη} έκπληξη — αυτό είναι παγίδα. Όμως το SUE
\emph{αυτοσυσχετίζεται}: μια θετική έκπληξη φέτος συνδέεται με θετική και το επόμενο τρίμηνο
(Bernard--Thomas). Το μετράμε ως απλή συσχέτιση:
\[ \rho_{\text{SUE}}=\mathrm{corr}\big(\mathrm{SUE}_{t},\,\mathrm{SUE}_{t+1}\big). \]
\ex{Στη δοκιμή βγήκε $\rho_{\text{SUE}}\approx0{,}315$ ενώ είχαμε βάλει $0{,}30$ — σωστή
ανάκτηση. Αυτό δίνει ένα μικρό, αμυνόμενο προγνωστικό στοιχείο για το επόμενο τρίμηνο.}

%==================================================================
\section{Σύνοψη συμβόλων}

\begin{center}
\renewcommand{\arraystretch}{1.35}
\begin{tabular}{@{}l l@{}}
\hline
\textbf{Σύμβολο} & \textbf{Τι σημαίνει}\\
\hline
$E_{i,t}$ & κέρδη ανά μετοχή της εταιρείας $i$ στο τρίμηνο $t$\\
$\sigma$ & τυπική απόκλιση («πόσο κουνιέται»)\\
$z$ & τυποποιημένη τιμή (z-score), κοινή κλίμακα\\
$\mathrm{SUE}$ & μέγεθος της έκπληξης κερδών $=(E_t-E_{t-4})/\sigma$\\
$\mathrm{AR}_{i,t}$ & abnormal απόδοση (αφαιρεμένη η αγορά)\\
$\mathrm{IC}=\rho$ & συσχέτιση σήματος με μελλοντική απόδοση\\
$b(k)$ & ολίσθηση ανά μονάδα SUE, τη μέρα $k$\\
$B(h)$ & σωρευτική ολίσθηση αν κρατήσεις $h$ μέρες\\
$\tau,\ t_{1/2}$ & χρόνος φθοράς και μισή ζωή $=\tau\ln 2$\\
$\mathrm{IR}$ & απόδοση ανά ρίσκο, ετήσια\\
$t_{\text{NW}}$ & $t$-στατιστική διορθωμένη (Newey--West)\\
$\rho_{\text{SUE}}$ & επιμονή της έκπληξης (forward στοιχείο)\\
\hline
\end{tabular}
\end{center}

\bigskip
\key{Η μεγάλη εικόνα: το \textbf{SUE} μετράει την έκπληξη· το \textbf{$b(k)$/half-life} λέει
\emph{πόσο} κρατάει· η \textbf{Newey--West $t$ \& IR} λένε \emph{αν είναι αληθινό}· και η
\textbf{durability} λέει \emph{αν αντέχει στον χρόνο}. Όλα από δωρεάν δεδομένα κερδών — χωρίς
προβλέψεις αναλυτών.}

\end{document}
